\documentclass[12pt]{article}
\usepackage[utf8]{inputenc}
\usepackage{amsmath, amssymb, amsthm}
\usepackage{geometry}
\usepackage{xcolor}
\usepackage[most]{tcolorbox}
\usepackage{enumitem}
\usepackage{fancyhdr}

% Page setup
\geometry{margin=1in}
\setlength{\headheight}{14.5pt}
\pagestyle{fancy}
\fancyhf{}
\rhead{AMC12 Problem Analysis}
\lhead{\today}
\cfoot{\thepage}

% Color definitions
\definecolor{problemconfig}{RGB}{230, 242, 255}    % Light blue
\definecolor{strategic}{RGB}{255, 230, 242}       % Light pink
\definecolor{tactical}{RGB}{242, 255, 230}        % Light green
\definecolor{techniques}{RGB}{255, 242, 230}      % Light yellow
\definecolor{verification}{RGB}{255, 230, 230}    % Light red
\definecolor{solution}{RGB}{230, 255, 255}        % Light cyan
\definecolor{competition}{RGB}{242, 242, 255}     % Light purple
\definecolor{gold}{RGB}{218, 165, 32}

% Box styles
\newtcolorbox{problemconfigbox}{
    colback=problemconfig,
    colframe=blue,
    title=Problem Configuration Analysis,
    fonttitle=\bfseries,
    arc=3pt,
    boxrule=1pt,
    breakable
}

\newtcolorbox{strategicbox}{
    colback=strategic,
    colframe=purple,
    title=Strategic Framework Application,
    fonttitle=\bfseries,
    arc=3pt,
    boxrule=1pt,
    breakable
}

\newtcolorbox{tacticalbox}{
    colback=tactical,
    colframe=green,
    title=Tactical Methodology Selection,
    fonttitle=\bfseries,
    arc=3pt,
    boxrule=1pt,
    breakable
}

\newtcolorbox{techniquesbox}{
    colback=techniques,
    colframe=orange,
    title=Conversion Techniques Application,
    fonttitle=\bfseries,
    arc=3pt,
    boxrule=1pt,
    breakable
}

\newtcolorbox{verificationbox}{
    colback=verification,
    colframe=red,
    title=Verification Strategy Design,
    fonttitle=\bfseries,
    arc=3pt,
    boxrule=1pt,
    breakable
}

\newtcolorbox{solutionbox}{
    colback=solution,
    colframe=cyan,
    title=Solution Roadmap,
    fonttitle=\bfseries,
    arc=3pt,
    boxrule=1pt,
    breakable
}

\newtcolorbox{competitionbox}{
    colback=competition,
    colframe=violet,
    title=Competition Strategy Integration,
    fonttitle=\bfseries,
    arc=3pt,
    boxrule=1pt,
    breakable
}

% Highlight boxes
\newtcolorbox{keyinsight}{
    colback=gold!10,
    colframe=gold,
    title=Key Insight,
    fonttitle=\bfseries,
    arc=3pt,
    boxrule=1.5pt,
    breakable
}

\newtcolorbox{warningbox}{
    colback=red!10,
    colframe=red,
    title=Warning,
    fonttitle=\bfseries,
    arc=3pt,
    boxrule=1.5pt,
    breakable
}

\newtcolorbox{protip}{
    colback=green!10,
    colframe=green!70!black,
    title=Pro Tip,
    fonttitle=\bfseries,
    arc=3pt,
    boxrule=1.5pt,
    breakable
}

\title{\textbf{AMC 12A 2024 Problem 21 Analysis}\\
\large{Comprehensive Strategic Framework Application}}
\author{Problem Analysis Framework v2.1}
\date{\today}

\begin{document}

\maketitle

\section*{Problem Statement}

Suppose that $a_1 = 2$ and the sequence $(a_n)$ satisfies the recurrence relation

\[\frac{a_n -1}{n-1}=\frac{a_{n-1}+1}{(n-1)+1}\]

for all $n \geq 2$. What is the greatest integer less than or equal to

\[\sum^{100}_{n=1} a_n^2?\]

\begin{align*}
\textbf{(A) } 338{,}550 \qquad \textbf{(B) } 338{,}551 \qquad \textbf{(C) } 338{,}552 \qquad \textbf{(D) } 338{,}553 \qquad \textbf{(E) } 338{,}554
\end{align*}

\section{Problem Configuration Analysis}

\begin{problemconfigbox}

\subsection*{A. Problem Classification}

\textbf{Problem Type:} To Find (computational problem with sequences and series)

\textbf{Difficulty Band:} Late-band (Problem 21 of 25)
\begin{itemize}
    \item Expected time: 5-7 minutes
    \item Difficulty level: High (84\% position)
    \item Requires recurrence relation solution
    \item Involves summation and estimation
\end{itemize}

\textbf{Competition Context:} AMC 12A (75 minutes, 25 problems, multiple choice)

\textbf{Mathematical Domain:} 
\begin{itemize}
    \item Primary: Sequences and Series (recurrence relations)
    \item Secondary: Algebra (manipulation and simplification)
    \item Tertiary: Summation Formulas (sum of squares, telescoping)
    \item Quaternary: Estimation (bounding infinite series)
\end{itemize}

\subsection*{B. Problem Structure Identification}

\textbf{Given Information:}
\begin{itemize}
    \item Initial condition: $a_1 = 2$
    \item Recurrence: $\frac{a_n -1}{n-1}=\frac{a_{n-1}+1}{n}$ for $n \geq 2$
    \item Goal: Find $\lfloor \sum_{n=1}^{100} a_n^2 \rfloor$
    \item Answer choices differ by only 1
\end{itemize}

\textbf{Constraints:}
\begin{itemize}
    \item Must solve recurrence to find explicit formula
    \item Must evaluate sum to 100 terms
    \item Need floor function (greatest integer)
    \item Answer choices very close together (requires precision)
\end{itemize}

\textbf{Objective:} Find the greatest integer $\leq \sum_{n=1}^{100} a_n^2$.

\textbf{Hidden Assumptions:}
\begin{itemize}
    \item The recurrence can be solved explicitly
    \item The sum will not be an integer (otherwise floor irrelevant)
    \item Small terms matter (answer choices differ by 1)
    \item Estimation techniques needed for $\sum \frac{1}{n^2}$
\end{itemize}

\textbf{Answer Choice Leverage:}
\begin{itemize}
    \item All choices in narrow range: 338,550-338,554
    \item Sequential: differ by exactly 1
    \item Large values suggest sum dominated by large terms
    \item Middle choice (C) might be trap if calculation error
    \item Need precise calculation to distinguish
\end{itemize}

\subsection*{C. Strategic Orientation}

\textbf{Hypothesis:} The recurrence transforms into a simpler form via substitution.

\textbf{Conclusion:} Solve for $a_n$, compute $a_n^2$, sum, and apply floor function.

\textbf{Notation:}
\begin{itemize}
    \item Let $b_n = n(a_n - 1)$ or $b_n = na_n$ (clever substitution)
    \item Use standard summation formulas
    \item Denote $S = \sum_{n=1}^{100} a_n^2$
\end{itemize}

\textbf{Intuitive Approach:} 
\begin{itemize}
    \item Transform recurrence via substitution
    \item Use telescoping to find explicit formula
    \item Recognize $a_n = n + \frac{1}{n}$ pattern
    \item Split sum into manageable parts
    \item Estimate $\sum \frac{1}{n^2}$ term
\end{itemize}

\textbf{Cross-Domain Potential:}
\begin{itemize}
    \item Algebraic: Manipulation and substitution
    \item Analysis: Series estimation and bounds
    \item Number theory: Floor function application
\end{itemize}

\end{problemconfigbox}

\begin{keyinsight}
\textbf{Critical Discovery:} The recurrence simplifies beautifully via substitution.

\textbf{Method 1 - Substitution $b_n = n(a_n - 1)$:}

Starting from: $\frac{a_n -1}{n-1}=\frac{a_{n-1}+1}{n}$

Cross-multiply: $n(a_n - 1) = (n-1)(a_{n-1} + 1)$

Let $b_n = n(a_n - 1)$, then:
\[b_n = (n-1)(a_{n-1} - 1) + 2(n-1) = b_{n-1} + 2(n-1)\]

This telescopes to: $b_n = b_1 + 2(1 + 2 + \cdots + (n-1)) = 1 + n(n-1) = n^2 - n + 1$

Therefore: $a_n = \frac{b_n + n}{n} = \frac{n^2 + 1}{n} = n + \frac{1}{n}$

\textbf{Method 2 - Substitution $b_n = na_n$:}

Transform to: $na_n - (n-1)a_{n-1} = 2n - 1$

Let $b_n = na_n$, then: $b_n - b_{n-1} = 2n - 1$

Telescoping: $b_n = b_1 + \sum_{k=2}^{n}(2k-1) = 2 + (3 + 5 + \cdots + (2n-1)) = 1 + n^2$

Therefore: $a_n = \frac{1 + n^2}{n} = n + \frac{1}{n}$

\textbf{Key Result:} $a_n = n + \frac{1}{n}$, so $a_n^2 = n^2 + 2 + \frac{1}{n^2}$
\end{keyinsight}

\section{Strategic Framework Application}

\begin{strategicbox}

\subsection*{A. Psychological Strategies}

\textbf{Mental Toughness:}
\begin{itemize}
    \item Problem 21 requires algebraic persistence
    \item Multiple steps: solve recurrence, sum, estimate
    \item Don't panic at complex recurrence
    \item Trust that clever substitution exists
\end{itemize}

\textbf{Creativity Cultivation:}
\begin{itemize}
    \item Think: What substitution simplifies this?
    \item Consider: Multiply both sides by something useful
    \item Pattern recognition: Look for telescoping
    \item Guess-and-check: Compute first few terms
\end{itemize}

\textbf{Iterative Refinement:}
\begin{itemize}
    \item First: Try to isolate $a_n$
    \item Second: If messy, try substitution
    \item Third: Look for pattern in small cases
    \item Fourth: Verify formula works
\end{itemize}

\subsection*{B. Investigation Strategies}

\textbf{Startup Orientation:}
\begin{itemize}
    \item Understand: This is a recurrence relation problem
    \item Identify: Need explicit formula for $a_n$
    \item Recognize: Substitution likely needed
    \item Note: Answer choices very close (need precision)
\end{itemize}

\textbf{Get Your Hands Dirty:}
\begin{itemize}
    \item Compute: $a_1 = 2$
    \item Compute: $a_2 = \frac{5}{2}$ from recurrence
    \item Compute: $a_3 = \frac{10}{3}$, $a_4 = \frac{17}{4}$
    \item Pattern: $a_n = \frac{n^2+1}{n}$ becomes evident!
\end{itemize}

\textbf{Penultimate Step:}
\begin{itemize}
    \item Final step: Apply floor function to sum
    \item Before that: Estimate $\sum \frac{1}{n^2}$ term
    \item Before that: Sum $n^2$ and constant terms
    \item Before that: Find $a_n^2 = n^2 + 2 + \frac{1}{n^2}$
    \item Before that: Solve for $a_n = n + \frac{1}{n}$
\end{itemize}

\textbf{Wishful Thinking:}
\begin{itemize}
    \item Wish: If only recurrence was simple $a_n = a_{n-1} + c$
    \item Reality: More complex, but substitution helps
    \item Wish: If only $\sum \frac{1}{n^2}$ was known exactly
    \item Reality: Need to bound between 1 and 2
\end{itemize}

\textbf{Make It Easier:}
\begin{itemize}
    \item Simplify: Cross-multiply recurrence first
    \item Simplify: Use substitution to linearize
    \item Simplify: Recognize telescoping pattern
    \item Simplify: Split sum into three manageable parts
\end{itemize}

\textbf{Conjecture via Small Cases:}
\begin{itemize}
    \item $a_1 = 2 = 1 + 1$
    \item $a_2 = \frac{5}{2} = 2 + \frac{1}{2}$
    \item $a_3 = \frac{10}{3} = 3 + \frac{1}{3}$
    \item $a_4 = \frac{17}{4} = 4 + \frac{1}{4}$
    \item Pattern: $a_n = n + \frac{1}{n}$ $\checkmark$
\end{itemize}

\textbf{Visualization:}
\begin{itemize}
    \item Visualize: Telescoping sum as cancellation
    \item Think: $\sum \frac{1}{n^2}$ converges to $\frac{\pi^2}{6} \approx 1.645$
    \item Understand: Why answer is 338,551 not 338,552
\end{itemize}

\subsection*{C. Late-Band Strategic Playbook}

\textbf{Answer-Choice Leverage:}
\begin{itemize}
    \item Sequential answers: 338,550 through 338,554
    \item Difference of 1: Small fractional part matters
    \item Pattern: $338,550 + k$ where $k \in \{0,1,2,3,4\}$
    \item Strategy: Once we get 338,550 + (1 to 2), answer is (B)
\end{itemize}

\textbf{Make It Easier Application:}
\begin{itemize}
    \item Replace: Complex recurrence with explicit formula
    \item Replace: Sum of $a_n^2$ with sum of simpler terms
    \item Replace: Exact calculation of $\sum \frac{1}{n^2}$ with bounds
    \item Build: Use known formula for $\sum n^2$
\end{itemize}

\textbf{Penultimate Step Focus:}
\begin{itemize}
    \item Working backwards: Need $\lfloor S \rfloor = 338,551$
    \item Therefore: Need $338,551 < S < 338,552$
    \item Therefore: Need $338,550 + 1 < S < 338,550 + 2$
    \item Check: $\sum \frac{1}{n^2} \in (1, 2)$ gives this $\checkmark$
\end{itemize}

\subsection*{D. Argument Construction Strategy}

\textbf{Direct Approach:}
\begin{enumerate}
    \item Cross-multiply the recurrence relation
    \item Apply clever substitution to linearize
    \item Use telescoping to find explicit formula
    \item Square $a_n$ and split into three sums
    \item Evaluate using standard formulas and bounds
    \item Apply floor function
\end{enumerate}

\textbf{Key Lemma:} $1 < \sum_{n=1}^{100} \frac{1}{n^2} < 2$

\textit{Proof:} Lower bound is trivial ($\frac{1}{1^2} = 1$). For upper bound, use comparison:
\[\sum_{n=1}^{\infty} \frac{1}{n^2} = \frac{\pi^2}{6} < 2\]
or group terms and bound by geometric series.

\end{strategicbox}

\begin{warningbox}
\textbf{Common Mistakes:}
\begin{enumerate}
    \item \textbf{Algebraic errors in recurrence manipulation:} Losing track of signs or coefficients
    \item \textbf{Forgetting initial condition:} Not verifying $a_1 = 2$ works with formula
    \item \textbf{Wrong summation formula:} Using $\frac{n(n+1)}{2}$ instead of $\frac{n(n+1)(2n+1)}{6}$
    \item \textbf{Misapplying floor function:} Taking floor too early or incorrectly
    \item \textbf{Poor estimation:} Not bounding $\sum \frac{1}{n^2}$ correctly
    \item \textbf{Arithmetic errors:} Large numbers make calculation mistakes likely
\end{enumerate}
\end{warningbox}

\section{Tactical Methodology Selection}

\begin{tacticalbox}

\subsection*{A. Core Mathematical Tactics}

\textbf{Substitution:}
\begin{itemize}
    \item Applied: Transform $a_n$ to $b_n$ for simpler recurrence
    \item Two approaches: $b_n = n(a_n-1)$ or $b_n = na_n$
    \item Goal: Linearize the recurrence
    \item Success: Creates telescoping sum
\end{itemize}

\textbf{Telescoping:}
\begin{itemize}
    \item Recognition: $b_n - b_{n-1} = f(n)$ pattern
    \item Application: Sum both sides from 2 to $n$
    \item Result: Most terms cancel, leaving simple expression
    \item Efficiency: Avoids iterating 100 times
\end{itemize}

\textbf{Pattern Recognition:}
\begin{itemize}
    \item Strategy: Compute first few terms
    \item Observation: $a_n = \frac{n^2+1}{n}$
    \item Verification: Check formula satisfies recurrence
    \item Advantage: Direct path if substitution unclear
\end{itemize}

\textbf{Estimation:}
\begin{itemize}
    \item Target: $\sum_{n=1}^{100} \frac{1}{n^2}$
    \item Method: Bound between simple values
    \item Lower: Clearly $> 1$
    \item Upper: Use $< \frac{\pi^2}{6} < 2$ or telescoping comparison
\end{itemize}

\subsection*{B. Domain-Specific Tactics}

\textbf{Sequences and Series Tactics:}
\begin{itemize}
    \item \textbf{Recurrence Solving:} Substitution, characteristic equations
    \item \textbf{Telescoping Sums:} Key for finding explicit formulas
    \item \textbf{Standard Formulas:} $\sum n^2 = \frac{n(n+1)(2n+1)}{6}$
    \item \textbf{Series Bounds:} Comparison tests, grouping methods
\end{itemize}

\textbf{Algebra Tactics:}
\begin{itemize}
    \item \textbf{Cross-Multiplication:} Clear fractions
    \item \textbf{Expansion:} $(n + \frac{1}{n})^2 = n^2 + 2 + \frac{1}{n^2}$
    \item \textbf{Factoring:} Recognize patterns
\end{itemize}

\textbf{Estimation Tactics:}
\begin{itemize}
    \item \textbf{Bounding Series:} Upper and lower bounds
    \item \textbf{Known Constants:} $\frac{\pi^2}{6} \approx 1.645$
    \item \textbf{Comparison:} $\frac{1}{n^2} < \frac{1}{n(n-1)}$ for $n \geq 2$
\end{itemize}

\subsection*{C. Crossover Techniques}

\textbf{Algebra $\leftrightarrow$ Sequences:}
\begin{itemize}
    \item Translate: Recurrence relation to algebraic equation
    \item Manipulate: Using standard algebraic techniques
    \item Return: To sequence interpretation
\end{itemize}

\textbf{Discrete $\leftrightarrow$ Analysis:}
\begin{itemize}
    \item Discrete: Finite sum to 100 terms
    \item Analysis: Use infinite series bounds
    \item Connection: Finite sum $<$ infinite sum
\end{itemize}

\end{tacticalbox}

\begin{protip}
\textbf{Standard Formula Quick Reference:}

Always have these memorized for AMC 12:

\begin{align*}
\sum_{n=1}^{N} n &= \frac{N(N+1)}{2} \\
\sum_{n=1}^{N} n^2 &= \frac{N(N+1)(2N+1)}{6} \\
\sum_{n=1}^{N} n^3 &= \left(\frac{N(N+1)}{2}\right)^2 \\
\sum_{n=1}^{\infty} \frac{1}{n^2} &= \frac{\pi^2}{6} \approx 1.645
\end{align*}

\textbf{Quick Calculation for $N=100$:}
\[\sum_{n=1}^{100} n^2 = \frac{100 \cdot 101 \cdot 201}{6} = \frac{2{,}030{,}100}{6} = 338{,}350\]

\textbf{Bounding $\sum \frac{1}{n^2}$:}

Lower bound: Obviously $> 1$ (first term alone)

Upper bound (telescoping method):
\[\frac{1}{n^2} < \frac{1}{n(n-1)} = \frac{1}{n-1} - \frac{1}{n}\]
\[\sum_{n=2}^{100} \frac{1}{n^2} < \sum_{n=2}^{100} \left(\frac{1}{n-1} - \frac{1}{n}\right) = 1 - \frac{1}{100} < 1\]
\[\sum_{n=1}^{100} \frac{1}{n^2} < 1 + 1 = 2\]
\end{protip}

\section{Conversion Techniques Application}

\begin{techniquesbox}

\subsection*{A. Recognition Index}

\textbf{Primary Technique: Substitution to Linearize}
\begin{itemize}
    \item Recognition trigger: Non-linear recurrence
    \item Strategy: Find substitution that creates linear form
    \item Candidates: $b_n = n(a_n - 1)$ or $b_n = na_n$
    \item Result: Telescoping sum emerges
\end{itemize}

\textbf{Secondary Technique: Pattern Recognition}
\begin{itemize}
    \item Recognition trigger: Small cases compute easily
    \item Strategy: Calculate $a_1, a_2, a_3, a_4$
    \item Pattern: $a_n = n + \frac{1}{n}$ becomes obvious
    \item Advantage: Bypasses substitution if pattern seen quickly
\end{itemize}

\textbf{Tertiary Technique: Series Decomposition}
\begin{itemize}
    \item Recognition trigger: Sum of $a_n^2$ with known form
    \item Strategy: Split into three separate sums
    \item Application: $\sum(n^2 + 2 + \frac{1}{n^2})$ separates cleanly
    \item Benefit: Each sum uses different technique
\end{itemize}

\subsection*{B. Technique Selection}

\textbf{Why Substitution Method?}
\begin{itemize}
    \item Systematic: Works even without pattern recognition
    \item Rigorous: Proves formula rather than guessing
    \item Generalizable: Technique applies to similar problems
    \item Elegant: Creates beautiful telescoping
\end{itemize}

\textbf{Why Pattern Recognition as Alternative?}
\begin{itemize}
    \item Fast: Immediate if pattern obvious
    \item Intuitive: Natural for visual thinkers
    \item Competition-friendly: Saves time under pressure
    \item Risk: Must verify formula works
\end{itemize}

\textbf{Why Series Estimation?}
\begin{itemize}
    \item Necessary: Can't compute $\sum \frac{1}{n^2}$ exactly easily
    \item Efficient: Bounds suffice for floor function
    \item Standard: Common competition technique
    \item Precise: Gives tight enough bounds
\end{itemize}

\subsection*{C. Integration Strategy}

\textbf{Phase 1: Solve Recurrence}
\begin{itemize}
    \item Cross-multiply original recurrence
    \item Apply substitution $b_n = n(a_n-1)$ or $b_n = na_n$
    \item Use telescoping to find $b_n$
    \item Extract $a_n = n + \frac{1}{n}$
\end{itemize}

\textbf{Phase 2: Square and Decompose}
\begin{itemize}
    \item Calculate $a_n^2 = n^2 + 2 + \frac{1}{n^2}$
    \item Write sum: $\sum a_n^2 = \sum n^2 + \sum 2 + \sum \frac{1}{n^2}$
    \item Recognize standard formulas apply
\end{itemize}

\textbf{Phase 3: Evaluate and Estimate}
\begin{itemize}
    \item Compute $\sum_{n=1}^{100} n^2 = 338,350$
    \item Compute $\sum_{n=1}^{100} 2 = 200$
    \item Estimate $1 < \sum_{n=1}^{100} \frac{1}{n^2} < 2$
    \item Sum: $338,550 + (1 \text{ to } 2)$
\end{itemize}

\textbf{Phase 4: Apply Floor Function}
\begin{itemize}
    \item Range: $S \in (338,551, 338,552)$
    \item Floor: $\lfloor S \rfloor = 338,551$
    \item Answer: (B)
\end{itemize}

\subsection*{D. Conflict Resolution}

\textbf{Potential Conflict 1: Which Substitution?}
\begin{itemize}
    \item Conflict: $b_n = n(a_n-1)$ vs. $b_n = na_n$
    \item Resolution: Both work! Choose whichever you see first
    \item Lesson: Multiple valid approaches exist
\end{itemize}

\textbf{Potential Conflict 2: Exact vs. Approximate}
\begin{itemize}
    \item Conflict: Need exact value of $\sum \frac{1}{n^2}$?
    \item Resolution: No! Bounds suffice for floor function
    \item Lesson: Don't overwork the problem
\end{itemize}

\textbf{Potential Conflict 3: Answer Choice Ambiguity}
\begin{itemize}
    \item Conflict: Sum could be near 338,551.5
    \item Resolution: Better bounds show it's in (338,551, 338,552)
    \item Lesson: Precision matters when choices are close
\end{itemize}

\end{techniquesbox}

\section{Verification Strategy Design}

\begin{verificationbox}

\subsection*{A. Verification Level Selection}

\textbf{Selected Level: Level 5 - Targeted Deep Check (2-3 minutes)}

\textbf{Decision Tree Analysis:}
\begin{itemize}
    \item Problem difficulty: Late-band (P21) $\to$ High stakes
    \item Personal confidence: Moderate (many calculation steps)
    \item Time remaining: Assume 12-15 minutes left $\to$ Can afford thorough check
    \item Answer type: Close multiple choices $\to$ Need precision
    \item Recommendation: Level 5 (Targeted Deep Check)
\end{itemize}

\subsection*{B. Specialized Verification Methods}

\textbf{Formula Verification:}

\textbf{Check 1: Does $a_n = n + \frac{1}{n}$ satisfy initial condition?}
\[a_1 = 1 + \frac{1}{1} = 2 \checkmark\]

\textbf{Check 2: Does formula satisfy recurrence?}
\begin{align*}
\frac{a_n - 1}{n-1} &= \frac{n + \frac{1}{n} - 1}{n-1} = \frac{n - 1 + \frac{1}{n}}{n-1} = 1 + \frac{1}{n(n-1)} \\
\frac{a_{n-1} + 1}{n} &= \frac{(n-1) + \frac{1}{n-1} + 1}{n} = \frac{n + \frac{1}{n-1}}{n} = 1 + \frac{1}{n(n-1)} \checkmark
\end{align*}

\textbf{Check 3: Verify pattern with computed values}
\begin{align*}
a_2 &= 2 + \frac{1}{2} = \frac{5}{2} \checkmark \\
a_3 &= 3 + \frac{1}{3} = \frac{10}{3} \checkmark \\
a_4 &= 4 + \frac{1}{4} = \frac{17}{4} \checkmark
\end{align*}

\textbf{Summation Verification:}

\textbf{Check 4: Sum of Squares Formula}
\begin{align*}
\sum_{n=1}^{100} n^2 &= \frac{100 \cdot 101 \cdot 201}{6} \\
&= \frac{2{,}030{,}100}{6} \\
&= 338{,}350 \checkmark
\end{align*}

\textbf{Check 5: Constant Sum}
\[\sum_{n=1}^{100} 2 = 2 \cdot 100 = 200 \checkmark\]

\textbf{Check 6: Bounds on $\sum \frac{1}{n^2}$}

Lower bound:
\[\sum_{n=1}^{100} \frac{1}{n^2} \geq \frac{1}{1^2} = 1 \checkmark\]

Upper bound (telescoping):
\begin{align*}
\sum_{n=1}^{100} \frac{1}{n^2} &= 1 + \sum_{n=2}^{100} \frac{1}{n^2} \\
&< 1 + \sum_{n=2}^{100} \frac{1}{n(n-1)} \\
&= 1 + \sum_{n=2}^{100} \left(\frac{1}{n-1} - \frac{1}{n}\right) \\
&= 1 + \left(1 - \frac{1}{100}\right) \\
&< 2 \checkmark
\end{align*}

\textbf{Check 7: Total Sum Range}
\begin{align*}
\sum a_n^2 &= 338{,}350 + 200 + \sum \frac{1}{n^2} \\
&= 338{,}550 + \sum \frac{1}{n^2} \\
&\in (338{,}550 + 1, 338{,}550 + 2) \\
&= (338{,}551, 338{,}552) \checkmark
\end{align*}

\textbf{Check 8: Floor Function Application}
\[\lfloor S \rfloor = 338{,}551 \text{ when } S \in (338{,}551, 338{,}552) \checkmark\]

\textbf{Answer Choice Validation:}

\textbf{Check 9: Does Answer Match?}
\begin{itemize}
    \item Our answer: 338,551
    \item Choice (B): 338,551 $\checkmark$
\end{itemize}

\textbf{Check 10: Why Not Adjacent Choices?}
\begin{itemize}
    \item (A) 338,550: Would need $\sum \frac{1}{n^2} < 1$ (impossible)
    \item (C) 338,552: Would need $\sum \frac{1}{n^2} \geq 2$ (false)
    \item Both ruled out $\checkmark$
\end{itemize}

\subsection*{C. Confidence Calibration}

\textbf{Pre-Verification Confidence: 80\%}
\begin{itemize}
    \item Strong: Clear derivation path
    \item Strong: Formula checks out for small cases
    \item Concern: Arithmetic errors in large numbers
    \item Concern: Bounds on series might be wrong
\end{itemize}

\textbf{Post-Verification Confidence: 97\%}
\begin{itemize}
    \item Confirmed: Formula satisfies recurrence
    \item Confirmed: All arithmetic correct
    \item Confirmed: Bounds are tight
    \item Confirmed: Answer matches choice
    \item Remaining 3\%: Minimal risk of oversight
\end{itemize}

\textbf{Final Decision: SELECT ANSWER (B)}

\end{verificationbox}

\section{Solution Roadmap}

\begin{solutionbox}

\subsection*{Complete Solution Path}

\textbf{Phase 1: Solve the Recurrence (2-3 minutes)}

\textbf{Step 1.1: Cross-Multiply the Recurrence}

Given: $\frac{a_n -1}{n-1}=\frac{a_{n-1}+1}{n}$

Cross-multiply:
\[n(a_n - 1) = (n-1)(a_{n-1} + 1)\]

\textbf{Step 1.2: Apply Substitution}

Let $b_n = n(a_n - 1)$. Then:
\begin{align*}
b_n &= (n-1)(a_{n-1} + 1) \\
&= (n-1)(a_{n-1} - 1) + 2(n-1) \\
&= b_{n-1} + 2(n-1)
\end{align*}

This is a simple recurrence: $b_n - b_{n-1} = 2(n-1)$

\textbf{Step 1.3: Use Telescoping}

Summing from $n=2$ to $n=N$:
\begin{align*}
b_N - b_1 &= \sum_{k=2}^{N} 2(k-1) \\
&= 2\sum_{k=1}^{N-1} k \\
&= 2 \cdot \frac{(N-1)N}{2} \\
&= N(N-1)
\end{align*}

Since $b_1 = 1(a_1 - 1) = 1(2-1) = 1$:
\[b_N = 1 + N(N-1) = N^2 - N + 1\]

\textbf{Step 1.4: Extract $a_n$}
\begin{align*}
n(a_n - 1) &= n^2 - n + 1 \\
a_n - 1 &= n - 1 + \frac{1}{n} \\
a_n &= n + \frac{1}{n}
\end{align*}

\textbf{Checkpoint 1:} Have I correctly solved for $a_n$?

\textbf{Phase 2: Square and Decompose (1 minute)}

\textbf{Step 2.1: Calculate $a_n^2$}
\begin{align*}
a_n^2 &= \left(n + \frac{1}{n}\right)^2 \\
&= n^2 + 2 \cdot n \cdot \frac{1}{n} + \frac{1}{n^2} \\
&= n^2 + 2 + \frac{1}{n^2}
\end{align*}

\textbf{Step 2.2: Write Sum as Three Parts}
\begin{align*}
\sum_{n=1}^{100} a_n^2 &= \sum_{n=1}^{100} \left(n^2 + 2 + \frac{1}{n^2}\right) \\
&= \sum_{n=1}^{100} n^2 + \sum_{n=1}^{100} 2 + \sum_{n=1}^{100} \frac{1}{n^2}
\end{align*}

\textbf{Checkpoint 2:} Have I correctly expanded $a_n^2$ and separated the sum?

\textbf{Phase 3: Evaluate Each Sum (2-3 minutes)}

\textbf{Step 3.1: Sum of Squares}
\begin{align*}
\sum_{n=1}^{100} n^2 &= \frac{100 \cdot 101 \cdot 201}{6} \\
&= \frac{100 \cdot 101 \cdot 201}{6} \\
&= \frac{2{,}030{,}100}{6} \\
&= 338{,}350
\end{align*}

\textbf{Step 3.2: Sum of Constants}
\[\sum_{n=1}^{100} 2 = 2 \cdot 100 = 200\]

\textbf{Step 3.3: Bound Sum of Reciprocal Squares}

Lower bound (trivial):
\[\sum_{n=1}^{100} \frac{1}{n^2} > \frac{1}{1^2} = 1\]

Upper bound (telescoping comparison):
\begin{align*}
\sum_{n=1}^{100} \frac{1}{n^2} &= 1 + \sum_{n=2}^{100} \frac{1}{n^2} \\
&< 1 + \sum_{n=2}^{100} \frac{1}{n(n-1)} \quad \text{(since } \frac{1}{n^2} < \frac{1}{n(n-1)} \text{ for } n \geq 2\text{)} \\
&= 1 + \sum_{n=2}^{100} \left(\frac{1}{n-1} - \frac{1}{n}\right) \\
&= 1 + \left(\frac{1}{1} - \frac{1}{2} + \frac{1}{2} - \frac{1}{3} + \cdots + \frac{1}{99} - \frac{1}{100}\right) \\
&= 1 + \left(1 - \frac{1}{100}\right) \\
&= 2 - \frac{1}{100} < 2
\end{align*}

Therefore: $1 < \sum_{n=1}^{100} \frac{1}{n^2} < 2$

\textbf{Checkpoint 3:} Have I correctly bounded the reciprocal sum?

\textbf{Phase 4: Final Calculation (30 seconds)}

\textbf{Step 4.1: Sum All Parts}
\begin{align*}
\sum_{n=1}^{100} a_n^2 &= 338{,}350 + 200 + \sum_{n=1}^{100} \frac{1}{n^2} \\
&= 338{,}550 + \sum_{n=1}^{100} \frac{1}{n^2}
\end{align*}

Since $1 < \sum_{n=1}^{100} \frac{1}{n^2} < 2$:
\[338{,}551 < \sum_{n=1}^{100} a_n^2 < 338{,}552\]

\textbf{Step 4.2: Apply Floor Function}
\[\lfloor \sum_{n=1}^{100} a_n^2 \rfloor = 338{,}551\]

\textbf{Step 4.3: Match to Answer}
\begin{itemize}
    \item Our answer: 338,551
    \item Choice: $\textbf{(B) } 338{,}551$
\end{itemize}

\textbf{Checkpoint 4:} Is my final answer correct?

\subsection*{Alternative Approaches}

\textbf{Alternative 1: Pattern Recognition Only}
\begin{itemize}
    \item Compute $a_1 = 2$, $a_2 = \frac{5}{2}$, $a_3 = \frac{10}{3}$, $a_4 = \frac{17}{4}$
    \item Recognize pattern: $a_n = \frac{n^2+1}{n}$
    \item Verify formula works for computed values
    \item Proceed directly to squaring and summing
    \item Faster if pattern seen quickly
\end{itemize}

\textbf{Alternative 2: Different Substitution}
\begin{itemize}
    \item Use $b_n = na_n$ instead of $b_n = n(a_n-1)$
    \item Transform: $na_n - (n-1)a_{n-1} = 2n - 1$
    \item Telescoping: $b_n = b_1 + \sum_{k=2}^{n}(2k-1) = 2 + (3+5+\cdots+(2n-1))$
    \item Result: $b_n = 1 + n^2$, so $a_n = \frac{1+n^2}{n} = n + \frac{1}{n}$
    \item Same result, different path
\end{itemize}

\textbf{Alternative 3: Use Known Value of $\frac{\pi^2}{6}$}
\begin{itemize}
    \item If you know $\sum_{n=1}^{\infty} \frac{1}{n^2} = \frac{\pi^2}{6} \approx 1.6449$
    \item Conclude $\sum_{n=1}^{100} \frac{1}{n^2} \approx 1.635$ (slightly less)
    \item Directly get $338{,}550 + 1.635 \approx 338{,}551.635$
    \item Floor gives 338,551
    \item Faster if constant memorized
\end{itemize}

\subsection*{Common Pitfalls}

\textbf{Pitfall 1: Algebraic Errors in Substitution}
\begin{itemize}
    \item Error: Incorrectly expanding $(n-1)(a_{n-1}+1)$
    \item Why wrong: Sign errors or missing terms
    \item Prevention: Carefully expand and check each step
\end{itemize}

\textbf{Pitfall 2: Wrong Summation Formula}
\begin{itemize}
    \item Common: Using $\sum n^2 = \frac{n(n+1)}{2}$ (wrong!)
    \item Correct: $\sum n^2 = \frac{n(n+1)(2n+1)}{6}$
    \item Prevention: Memorize standard formulas
\end{itemize}

\textbf{Pitfall 3: Poor Bounds on $\sum \frac{1}{n^2}$}
\begin{itemize}
    \item Error: Not bounding tightly enough
    \item Why wrong: Might not distinguish between answer choices
    \item Prevention: Use telescoping comparison to get $< 2$
\end{itemize}

\textbf{Pitfall 4: Arithmetic Errors}
\begin{itemize}
    \item Common: $\frac{2{,}030{,}100}{6} = 338{,}350$ miscalculated
    \item Prevention: Double-check division with large numbers
    \item Tip: $\frac{100 \cdot 101 \cdot 201}{6} = \frac{100 \cdot 101 \cdot 201}{6}$
\end{itemize}

\textbf{Pitfall 5: Floor Function Confusion}
\begin{itemize}
    \item Error: Rounding instead of flooring
    \item Why wrong: Floor always rounds down
    \item Example: $\lfloor 338{,}551.9 \rfloor = 338{,}551$, not 338,552
\end{itemize}

\end{solutionbox}

\section{Competition Strategy Integration}

\begin{competitionbox}

\subsection*{A. Time Management}

\textbf{Time Budget for Problem 21:}
\begin{itemize}
    \item Total allocated: 5-7 minutes (late-band problem)
    \item Phase 1 (Solve recurrence): 2-3 minutes
    \item Phase 2 (Square/decompose): 1 minute
    \item Phase 3 (Evaluate sums): 2-3 minutes
    \item Phase 4 (Final calculation): 30 seconds
    \item Phase 5 (Verification): 1-2 minutes
\end{itemize}

\textbf{Time Checkpoints:}
\begin{itemize}
    \item After 3 minutes: Should have $a_n = n + \frac{1}{n}$
    \item After 5 minutes: Should have sum decomposed and evaluated
    \item After 7 minutes: Should have final answer
    \item After 8 minutes: If not done, make educated guess
\end{itemize}

\subsection*{B. Strategic Positioning}

\textbf{When to Attempt:}
\begin{itemize}
    \item Strong algebra: Attempt after P1-20
    \item Weak sequences: Consider skipping for P22-25
    \item If $>$12 minutes remaining: Attempt
    \item If familiar with recurrences: High success probability
\end{itemize}

\textbf{Problem Sequencing:}
\begin{itemize}
    \item This is P21 - beginning of late band
    \item More accessible than P22-25 for some students
    \item Consider: Algebraic vs. geometric/combinatorial preference
    \item Strategy: Attempt if algebra is strength
\end{itemize}

\subsection*{C. Risk Assessment}

\textbf{Success Probability:}
\begin{itemize}
    \item With strong algebra: 70-80\%
    \item With recurrence experience: 75-85\%
    \item Standard preparation: 50-60\%
    \item If pattern seen quickly: 80-90\%
\end{itemize}

\textbf{Risk Level: MEDIUM}
\begin{itemize}
    \item Difficulty: High (P21/25)
    \item Complexity: Multiple steps but systematic
    \item Time cost: 5-7 minutes
    \item Payoff: Same 6 points as easier problems
    \item Advantage: More systematic than P23-24
\end{itemize}

\textbf{Risk Mitigation:}
\begin{itemize}
    \item Quick check: Can I compute $a_2, a_3$? (If yes: proceed)
    \item Quick check: Do I see pattern in small cases? (If yes: fast path)
    \item Quick check: Do I know $\sum n^2$ formula? (If no: risky)
    \item Backup: Answer choices differ by 1, so precision matters
\end{itemize}

\subsection*{D. Abandonment Criteria}

\textbf{When to Abandon:}

\textbf{Scenario 1: After 3 minutes, can't solve recurrence}
\begin{itemize}
    \item Action: Try pattern recognition with small cases
    \item If still stuck: Skip to next problem
    \item Reason: Fundamental blockage unlikely to resolve
\end{itemize}

\textbf{Scenario 2: After 6 minutes, arithmetic errors mounting}
\begin{itemize}
    \item Action: Use answer choice structure to guess
    \item Reasoning: Middle choices (B,C,D) most likely
    \item Best guess: (B) based on bound analysis
\end{itemize}

\textbf{Scenario 3: Less than 10 minutes remaining in exam}
\begin{itemize}
    \item Action: Strategic guess immediately
    \item Reason: Not enough time for full solution
\end{itemize}

\textbf{Strategic Guessing (if must abandon):}
\begin{itemize}
    \item Observe: Answers are 338,550 through 338,554
    \item Pattern: Sequential, differ by 1
    \item Key insight: $\sum \frac{1}{n^2}$ is between 1 and 2
    \item Therefore: Answer is 338,550 + 1 = 338,551 or 338,550 + 2 = 338,552
    \item Best guess: (B) 338,551 (sum is closer to 1.6 than to 2)
\end{itemize}

\subsection*{E. Answer Recording}

\textbf{Recording Protocol:}
\begin{itemize}
    \item Write answer clearly: B
    \item Double-check bubble: (B) is filled
    \item Cross-reference: Problem 21 = bubble 21
    \item Watch for: Large numbers, easy to misread
\end{itemize}

\subsection*{F. Post-Problem Strategy}

\textbf{If Solved Successfully:}
\begin{itemize}
    \item Confidence: Algebra skills validated
    \item Time check: How much remaining?
    \item Next: Tackle remaining late-band problems
\end{itemize}

\textbf{If Abandoned:}
\begin{itemize}
    \item No regret: P21 requires specific techniques
    \item Focus: Maximize points on achievable problems
    \item Learn: Review telescoping and substitution methods later
\end{itemize}

\end{competitionbox}

\section{Summary and Key Takeaways}

\begin{keyinsight}
\textbf{Critical Insights for This Problem:}

\begin{enumerate}
    \item The recurrence relation simplifies beautifully via substitution. Either $b_n = n(a_n-1)$ or $b_n = na_n$ works, creating a telescoping sum.
    
    \item The explicit formula $a_n = n + \frac{1}{n}$ can be found either by solving the recurrence or by computing small cases and recognizing the pattern.
    
    \item Squaring gives $a_n^2 = n^2 + 2 + \frac{1}{n^2}$, which decomposes into three manageable sums.
    
    \item The key to getting the right answer is bounding $\sum_{n=1}^{100} \frac{1}{n^2}$ correctly between 1 and 2. The telescoping comparison $\frac{1}{n^2} < \frac{1}{n(n-1)}$ is the elegant method.
    
    \item With $338,550 + (1 \text{ to } 2)$, the floor gives $338,551$. Answer choices differ by 1, so precision is essential.
\end{enumerate}
\end{keyinsight}

\begin{protip}
\textbf{General Strategy for Recurrence Problems:}

When faced with recurrence relations:

\begin{enumerate}
    \item \textbf{Try small cases first:} Compute $a_1, a_2, a_3, a_4$ and look for patterns. Often the pattern is obvious and you can guess-and-verify.
    
    \item \textbf{Look for simplifying substitutions:} If pattern unclear, try:
    \begin{itemize}
        \item Multiplying/dividing by $n$
        \item Adding/subtracting related terms
        \item Creating combinations like $b_n = n \cdot a_n$ or $b_n = a_n - f(n)$
    \end{itemize}
    
    \item \textbf{Aim for telescoping:} The ideal transformed recurrence has form $b_n - b_{n-1} = f(n)$, which telescopes cleanly.
    
    \item \textbf{Verify your formula:} Always check that your explicit formula satisfies both the recurrence and initial condition.
    
    \item \textbf{Use standard formulas:} Memorize $\sum n$, $\sum n^2$, $\sum n^3$ and know how to bound simple series.
\end{enumerate}

For this problem, recognizing that $b_n = n(a_n - 1)$ or $b_n = na_n$ creates a linear recurrence is the breakthrough that unlocks everything.
\end{protip}

\section{Final Answer}

\begin{center}
\Large
\fbox{\textbf{Answer: (B) } $338{,}551$}
\end{center}

\textbf{Solution Summary:}
\begin{itemize}
    \item Transform recurrence via substitution: $b_n = n(a_n-1)$ gives $b_n - b_{n-1} = 2(n-1)$
    \item Telescope to find: $b_n = n^2 - n + 1$
    \item Extract explicit formula: $a_n = n + \frac{1}{n}$
    \item Square: $a_n^2 = n^2 + 2 + \frac{1}{n^2}$
    \item Sum parts:
    \begin{align*}
    \sum_{n=1}^{100} a_n^2 &= \sum_{n=1}^{100} n^2 + \sum_{n=1}^{100} 2 + \sum_{n=1}^{100} \frac{1}{n^2} \\
    &= 338{,}350 + 200 + (1 \text{ to } 2) \\
    &\in (338{,}551, 338{,}552)
    \end{align*}
    \item Apply floor: $\lfloor S \rfloor = 338{,}551$
\end{itemize}

\end{document}